\section{Introduction}
\label{sec:introduction}

The board exists to drive six cold-cathode gas-discharge indicator
tubes from a USB-C Power Delivery source: four IN-12B numeric tubes
forming a HH:MM digit display, and two IN-6 tubes forming a colon
separator between the hour and minute pairs. Each numeric tube is
paired one-to-one with a K155ID1 BCD decoder/driver, which
selects one of ten cathodes per digit under control of an MCU. Every
other subsystem on the board --- the 170~V anode supply, the 5~V rail
for the decoder/drivers, the 3.3~V rail for the MCU, and the USB-PD
front end that supplies all three --- exists only to put a stable,
correctly proportioned high voltage across these tubes and to switch
it through these driver chips. The tubes and drivers are therefore the
correct starting point for a design analysis: their datasheet ratings
set the requirements that every upstream converter stage must satisfy,
rather than the other way around.

This document derives those requirements first (\cref{sec:tubes,sec:bcd-driver}),
then works backward through the 170~V flyback converter that meets
them (\cref{sec:hv-components,sec:hv-converter}), the 5~V rail that
powers the decoder/drivers and the 3.3~V rail that powers the MCU
(\cref{sec:lv-components,sec:5v-converter,sec:3v3-converter}), and
finally the USB-PD input stage sized to supply all three downstream
converters.

\paragraph{Method.} Every parameter attributed to a named,
design-driving component --- a Nixie tube, the BCD decoder/driver, the
flyback controller, the flyback transformer --- is transcribed from a
datasheet on file in \texttt{docs/datasheets/} and cited by filename
and page. Where a datasheet gives a range, a tolerance, or a specific
test condition, that context is preserved rather than collapsed to a
single number, because the converter design in
\cref{sec:hv-converter} depends on knowing which figures are
guaranteed worst cases and which are typical or test-condition values.
\texttt{docs/datasheets/} is the boundary of that treatment: a part
gets a cited, traceable figure if and only if its datasheet lives
there, and everything else --- resistors, capacitors, standard diodes,
and any other component the design doesn't hinge on --- is sized to
ordinary engineering practice instead. Package and hand-soldering
constraints enter only where a single critical part is available in
one specific package (for example, the flyback transformer or the
load-switch dual MOSFET); otherwise the analysis stays at the
schematic level.

\paragraph{A note on the BCD driver substitution.} No datasheet for
the K155ID1 itself --- a Soviet-made functional clone of the Texas
Instruments SN74141 --- could be located. By explicit decision, this
document treats the K155ID1 as electrically equivalent to a generic
TI SN74141 and cites that datasheet in its place
(\cref{sec:bcd-driver}). This is a documented substitution, not a
verified characteristic of the part actually populated on the board,
and it is flagged again at the point where it matters most: the
driver's off-state output voltage rating.
